\begin{frame}{FAQ (1/2): $V$ 와 $Q$ 의 연결, 그리고 ``최적''}
  \begin{itemize}
    \item \textbf{Q. $V^\pi(s)$ 와 $Q^\pi(s,a)$ 는 서로 어떻게
      연결되나요?}
      \begin{itemize}
        \item $V^\pi(s) = \sum_a \pi(a|s)\,Q^\pi(s,a)$ --- 정책이
          각 행동을 고를 \textbf{확률로 $Q^\pi$ 를 가중평균}.
        \item 반대로 $Q^\pi(s,a) = R(s,a) + \gamma\,
          \mathbb{E}_{s'\sim P(\cdot|s,a)}[V^\pi(s')]$ --- ``이번
          행동만 $a$ 로 \textbf{강제로 고정}하고, 그 다음부터는
          $\pi$ 를 따른다''.
        \item 두 함수는 \textbf{서로를 정의하는 데 쓰이는, 동전의
          양면} 같은 관계.
      \end{itemize}
    \item \textbf{Q. ``최적'' 정책의 가치함수는 어떻게
      정의하나요?}
      \begin{itemize}
        \item 이번 절의 $V^\pi, Q^\pi$ 는 \textbf{주어진} 정책
          $\pi$ 를 따랐을 때의 가치를 잰다.
        \item 강화학습의 궁극 목표는 ``가능한 모든 정책 중 가치가
          가장 높은'' \kb{최적 정책 $\pi^*$}를 찾는 것인데, 그
          정의와 계산법(\kb{벨만 \textbf{최적}방정식} ---
          이번 절의 벨만 \textbf{기대}방정식과 한 글자 차이지만
          $\max$ 연산이 들어가 \textbf{의미가 크게 달라진다})는
          Week 4에서 바로 이어 다룬다.
      \end{itemize}
  \end{itemize}
\end{frame}
